\section{Conclusion}
\label{sec:conclusion}

This paper has given a technical account of Vibe Theory as it now stands: the model in
its current form, and thirty-six findings from a finite, reproducible testbed.

The model is one legible object. It treats reality as a single kind of entity, the
vibe, a unit of experience carrying a ternary tone and related only by the note, on a
growing, random, hyperbolic, causal mesh, updated by a local ternary signed-majority
rule. There are no
continuous weights and no hidden state, and every continuous quantity of physics is an
emergent, coarse-grained description of the discrete substrate. The model is written and
run at a glance by a small constructor, so its description and its implementation are
the same object.

In one picture, the structure is a growing, random, hyperbolic network of three-valued
points that vote locally, a causal set carrying a network automaton. It is deliberately
not a regular grid or a regular tiling, which would single out a preferred frame and
break Lorentz invariance, and not a field on a smooth space, since there is none
underneath. The smooth quantities of physics, the geometry, the energy and flow of time,
and the quantum amplitude, are large-scale summaries of the web rather than values
stored on any point. That single object is what every finding above was measured on.

The findings, taken together, make two claims. Monism is computationally coherent and
expressive: geometry, time, the law, matter and the four field spins, mass, force,
gravity, the quantum pillars, the dark sector, and cosmology all come out of the one
mesh, with the gauge and graviton operators derived from the action and the electroweak
masses reproduced. And the model runs as one system: a single instantiation of the
present model yields a definite Lorentz-safe geometry, stable structured states, the
local bounded-below Hamiltonian, and the arrow of time. That is the difference between a
pile of demonstrations on convenient substrates and one model that works end-to-end.

Three of the framework's distinctive predictions make contact with data. The everpresent
cosmological constant matches the observed dark energy to order of magnitude, the right
value emerging as one over the square root of the spacetime volume rather than a tuned
constant. The absence of linear Lorentz violation is confirmed by gamma-ray-burst
timing, and it follows from the random sprinkling in a way a lattice could not survive.
And the swerve is a clean future target.

We have been equally precise about the limits. The Born rule is shown consistent, not
derived. The Einstein equation is linearized, not yet nonlinear. The Standard-Model
content beyond the electroweak masses, Hawking radiation, the large-$N$ free-energy
crossing, and a fully emergent spatial geometry from a growth rule are open. The
framework's foundational claim about experience is a posit, outside what any simulation
can decide.

What Vibe Theory offers, then, is a candidate ontology that has been written down as one
legible model, simulated rather than only argued, validated against known physics where
it can be, and confronted with data where its predictions are distinctive. Its physical
ladder is built and largely validated, its hardest dynamical question is measured, its
quantum question is sharpened to a precise open problem, two of its predictions already
meet observation, and its boundary is named honestly. The substrate is one kind of
thing, related to itself in different patterns, and a great deal of the world really does
come out of it.
